\chapter{Federated Learning with Differential Privacy}
\label{ch:federated_learning}

\section{Introduction to Federated Learning}
\label{sec:fl_introduction}

\subsection{Motivation}
\label{subsec:fl_motivation}

Centralized training in large-scale distributed systems presents several 
limitations:

\begin{itemize}
    \item \textbf{Data Volume}: Aggregating raw telemetry or operational data 
    from thousands of distributed nodes may exceed available network bandwidth 
    and storage capacity, creating bottlenecks and increasing latency in model 
    training \cite{li2020federated}.
    
    \item \textbf{Privacy Constraints}: In many industrial and regulatory contexts, 
    sensitive telemetry or Personally Identifiable Information (PII) must not 
    leave the local environment. Regulations such as the General Data Protection 
    Regulation (GDPR) \cite{gdpr2016regulation} impose legal obligations to limit 
    cross-boundary data transfer \cite{kairouz2021advances}.
    
    \item \textbf{Network Bandwidth Limitations}: Edge nodes often operate over 
    constrained or intermittent networks. Continuously streaming raw datasets is 
    infeasible, leading to delays or dropped data \cite{bonawitz2019towards}.
    
    \item \textbf{Security Risks}: Transmitting raw data increases the attack 
    surface for adversaries who may intercept or tamper with sensitive information 
    \cite{xu2021federated}.
\end{itemize}

Federated Learning (FL) addresses these challenges by enabling nodes to train 
locally while sharing only model updates or gradients, significantly reducing 
communication overhead and improving data privacy \cite{mcmahan2017communication}.

\subsection{Federated Learning Paradigm}
\label{subsec:fl_paradigm}

Federated Learning is a distributed machine learning paradigm in which a global 
model is collaboratively trained across multiple clients, each holding private 
local datasets \cite{mcmahan2017communication}. Let there be $K$ clients, each 
with dataset $\mathcal{D}_k$. The global objective function is:

\begin{equation}
    \min_{\theta} F(\theta) = \sum_{k=1}^{K} \frac{|\mathcal{D}_k|}{\sum_{j=1}^{K} |\mathcal{D}_j|} F_k(\theta),
\end{equation}

where $F_k(\theta)$ denotes the local loss at client $k$. Each federated learning 
round comprises:

\begin{enumerate}
    \item \textbf{Local Training}: Clients update model parameters $\theta$ using 
    their local private data.
    
    \item \textbf{Gradient or Weight Upload}: Clients transmit parameter updates 
    $\Delta \theta_k$ or gradients $\nabla F_k(\theta)$ to the central server.
    
    \item \textbf{Aggregation}: The server computes a global update by aggregating 
    received updates.
    
    \item \textbf{Broadcast}: The updated global model $\theta_{t+1}$ is distributed 
    back to all clients for the next round.
\end{enumerate}

This iterative process continues until convergence criteria, such as loss 
stabilization or a predefined number of rounds are met \cite{li2020federated}.

\section{Federated Learning Architecture}
\label{sec:fl_architecture}

\subsection{Participants: Edge Nodes and Central Coordinator}
\label{subsec:fl_participants}

The federated learning architecture comprises two participant roles.

\textbf{Edge Nodes (Clients).} Each edge node monitors a single microservice
instance and maintains a local copy of the LSTM-based autoencoder (LSTM-AE)
described in Chapter~\ref{ch:model_architecture}. During each FL round, the
node collects a window of local telemetry---CPU utilization, memory consumption,
request latency, and error counts---and trains the local model for a small
number of epochs. The resulting parameter update $\Delta\theta_k$ is compressed
using the adaptive serialization and compression pipeline described in
Chapter~\ref{ch:communication} before being transmitted to the coordinator via
gRPC. Between rounds, the local model continues to perform real-time inference
independently, ensuring that anomaly detection is not interrupted by network
partitions or coordinator downtime.

\textbf{Central Coordinator (Server).} The coordinator maintains the
authoritative global model $\theta_t$ and orchestrates each FL round. It (i)
registers edge nodes and tracks their participation state, (ii) receives and
buffers compressed model updates, (iii) verifies updates against a poisoning
detection pipeline (Section~\ref{sec:fl_fault_tolerance}), (iv) aggregates
valid updates using FedAvg, and (v) broadcasts the updated global model back to
all registered nodes. The coordinator persists its state---including the current
round number, the global model checkpoint, and the client registry---to a
durable database, allowing it to resume after a crash without losing progress
(see Chapter~\ref{ch:reliability}).

\subsection{Aggregation Strategies}
\label{subsec:fl_aggregation}

This framework adopts the \textit{Federated Averaging} (FedAvg) algorithm for 
global model aggregation, which is widely adopted due to its simplicity and 
communication efficiency \cite{mcmahan2017communication}.

FedAvg computes a weighted average of client updates according to their local dataset 
size. Formally, if $\theta_t^{(k)}$ is the updated model from client $k$ after 
local training, the server computes:

\begin{equation}
    \theta_{t+1} = \sum_{k=1}^{K} \frac{n_k}{\sum_{j=1}^{K} n_j} \theta_t^{(k)},
\end{equation}

where $n_k = |\mathcal{D}_k|$. This weighting ensures that clients with more 
data contribute proportionally to the global update. Under the assumptions of 
smooth loss functions, bounded gradients, and Independent and Identically 
Distributed (IID) or mildly non-IID data, FedAvg converges efficiently to a 
stationary point \cite{li2020federated, mcmahan2017communication}.

FedAvg’s properties make it particularly suitable for resource-constrained edge 
environments:

\begin{itemize}
    \item \textbf{Low Communication Overhead}: Only model weights or gradients 
    are transmitted, avoiding raw data movement.
    
    \item \textbf{Computational Simplicity}: Clients perform standard local 
    Stochastic Gradient Descent (SGD) updates without additional complex 
    computations.
    
    \item \textbf{Robustness to Moderate Heterogeneity}: Works effectively when 
    client data distributions vary slightly \cite{li2020federated}.
\end{itemize}

\subsection{Training and Evaluation in a Federated Setting}
\label{subsec:fl_training}

Each edge node performs local training using the following procedure.

\begin{enumerate}
    \item \textbf{Initialization.} At the start of a round, the node loads the
    latest global model parameters $\theta_t$ received from the coordinator.
    The current parameters are saved as $\theta_t^{\text{init}}$ so that the
    model update $\Delta\theta_k = \theta_k - \theta_t^{\text{init}}$ can be
    computed after training.

    \item \textbf{Local Training.} The node trains for $E=5$ local epochs on
    its private dataset using the Adam optimiser with a learning rate of
    $10^{-3}$ and a batch size of~64. The training objective is reconstruction
    loss (MSE) between the input window and the LSTM-AE output.

    \item \textbf{Gradient Compression.} After training, the parameter delta
    $\Delta\theta_k$ undergoes optional Top-K sparsification
    (Section~\ref{sec:compression}) with a default compression ratio of 10\%,
    retaining only the largest-magnitude updates and accumulating the remainder
    in a local error-feedback buffer.

    \item \textbf{Differential Privacy.} If DP is enabled, gradients are
    clipped to a norm of $C=1.0$ and Gaussian noise calibrated to the privacy
    budget $\varepsilon$ is added before transmission
    (Section~\ref{sec:differential_privacy}).

    \item \textbf{Transmission.} The compressed, optionally noised update is
    serialized and adaptively compressed at the byte level before being
    transmitted to the coordinator over a gRPC channel.
\end{enumerate}

\textbf{Communication frequency.} By default, nodes synchronize with the
coordinator once per round. The coordinator waits up to a configurable timeout
(default 120\,s) for at least $K_{\min}$ client updates before triggering
aggregation. Nodes that miss a round simply re-join in the next one.

\textbf{Evaluation.} After each aggregation, the global model can be evaluated
on a held-out validation partition. In the benchmark experiments
(Chapter~\ref{ch:evaluation}), the global model is evaluated after every round,
yielding F1, Precision, and Recall on a separate test set with injected
anomalies.

\subsection{Model Deployment to Edge Nodes}
\label{subsec:fl_deployment}

After each successful aggregation round, the updated global model is deployed
back to all registered edge nodes.

\textbf{Version management.} The coordinator maintains a monotonically
increasing round counter $t$. Each global model checkpoint is tagged with this
round number and persisted to durable storage. A configurable retention policy
keeps the last $N$ checkpoints, enabling rollback if a newly aggregated model
exhibits degraded detection quality.

\textbf{Rollout strategy.} Deployment follows a pull-based model: edge nodes
request the latest global model from the coordinator via the
\texttt{GetGlobalModel} gRPC endpoint. The coordinator responds with the
serialized and compressed model state along with the round number. This
approach avoids broadcast storms and allows nodes to update at their own pace.

\textbf{Offline nodes.} If a node was unreachable during a round, it is not
penalized. When the node reconnects, it re-registers with the coordinator and
downloads the latest global model. Its \texttt{last\_seen} timestamp is updated,
and it is enrolled in subsequent rounds automatically. Staleness tolerance is
configurable; updates from nodes whose local round number lags behind the global
round by more than a configured tolerance (default~2 rounds) are discarded
during aggregation.

\textbf{Hot-swapping.} The local model is swapped atomically: the node
continues inference with the previous model weights until the new state
dictionary has been fully loaded into memory. The swap operation involves a
single call to \texttt{load\_state\_dict()}, after which the model resumes
evaluation mode. Because torch operations are synchronous and the swap is
in-memory, no inference windows are dropped during the transition.

\section{Differential Privacy}
\label{sec:differential_privacy}

\subsection{Motivation for Differential Privacy}
\label{subsec:dp_motivation}

Even in federated learning, model updates can leak information about individual 
training samples:

\begin{itemize}
    \item \textbf{Membership Inference}: An adversary may determine whether a 
    specific sample was included in training by observing the changes in model 
    parameters \cite{shokri2017membership}.
    
    \item \textbf{Model Inversion}: Sensitive attributes can be partially 
    reconstructed from model weights or gradients \cite{fredrikson2015model}.
\end{itemize}

Differential privacy (DP) provides a formal framework to bound the influence of 
any single sample on the released updates. Formally, a randomized mechanism 
$\mathcal{M}$ satisfies $(\varepsilon, \delta)$-DP if for any pair of 
neighboring datasets $D, D'$ differing in a single sample, and for all measurable 
subsets $S$ of outputs:

\begin{equation}
    \Pr[\mathcal{M}(D) \in S] \leq e^\varepsilon \Pr[\mathcal{M}(D') \in S] + \delta.
\end{equation}

This ensures that an observer cannot confidently infer the presence or absence 
of any individual record from the mechanism’s output \cite{dwork2014algorithmic, 
abadi2016deep}.

\subsection{How Differential Privacy Works}
\label{subsec:dp_mechanism}

DP mechanisms applied in gradient-based learning typically involve the following 
steps:

\begin{enumerate}
    \item \textbf{Per-Sample Gradient Clipping}: For each sample $x_i$ in a 
    mini-batch of size $B$, the gradient $g_i = \nabla_\theta \mathcal{L}(\theta, x_i)$ 
    is clipped to a threshold $C$:

    \begin{equation}
        \bar{g}_i = g_i \cdot \min \left(1, \frac{C}{\| g_i \|_2} \right),
    \end{equation}

    limiting the influence of any single training example on the aggregate update.

    \item \textbf{Noise Addition}: The aggregated gradient over the batch is 
    perturbed with Gaussian noise:

    \begin{equation}
        \tilde{g} = \frac{1}{B} \sum_{i=1}^{B} \bar{g}_i + \mathcal{N}\left(0, \sigma^2 C^2 \mathbf{I}\right),
    \end{equation}

    where $\sigma$ is set according to the desired $(\varepsilon, \delta)$ privacy 
    budget and $\mathbf{I}$ is the identity matrix.

    \item \textbf{Privacy Budget Accounting}: Across $T$ rounds, the overall privacy 
    loss is composed using advanced composition \cite{dwork2014algorithmic} or 
    the moments accountant method \cite{abadi2016deep}, allowing tight estimation 
    of cumulative $\varepsilon$:

    \begin{equation}
        \varepsilon_\text{total} = f(\varepsilon_\text{per\_round}, T, \delta),
    \end{equation}

    where $f$ depends on the composition theorem applied.
    
    \item \textbf{Trade-offs}: Increasing the noise scale $\sigma$ improves 
    privacy by reducing the per-round $\varepsilon$, but it may degrade model 
    accuracy. Selecting appropriate values for the clipping norm $C$, the noise 
    scale $\sigma$, and the number of local training steps balances the privacy 
    budget $(\varepsilon, \delta)$ against utility \cite{geyer2017differentially}.
\end{enumerate}

\subsection{Differential Privacy in Federated Learning}
\label{subsec:dp_in_fl}

In federated learning, DP is implemented at the client level to ensure that updates 
shared with the central server do not leak individual sample information:

\begin{itemize}
    \item \textbf{Local Clipping}: Each client clips gradients per-sample before 
    aggregation to bound sensitivity:

    \begin{equation}
        \bar{g}_k = \frac{1}{|\mathcal{D}_k|} \sum_{i \in \mathcal{D}_k} 
        g_i \cdot \min \left(1, \frac{C}{\| g_i \|_2} \right),
    \end{equation}

    where $\mathcal{D}_k$ is the dataset of client $k$.

    \item \textbf{Noise Calibration}: Gaussian noise $\mathcal{N}(0, \sigma^2 C^2 \mathbf{I})$ 
    is added to each client’s aggregated update according to the specified 
    privacy budget.

    \item \textbf{Composition Across Rounds}: The cumulative privacy loss over 
    $T$ FL rounds is formally tracked, often using the moments accountant, to 
    ensure that $(\varepsilon, \delta)$-DP guarantees hold for the entire training 
    process.
\end{itemize}

These mechanisms ensure that even if the central server or a malicious participant 
observes all client updates, the influence of any individual training sample is 
provably limited, enabling privacy-preserving federated learning in edge and 
distributed environments.

\section{Fault Tolerance in Federated Learning}
\label{sec:fl_fault_tolerance}

\subsection{Adaptive Node Profiles}
\label{subsec:fl_node_profiles}

To accommodate the heterogeneity of edge deployments, the framework defines two
node profiles that govern fault-tolerance behavior.

\textbf{Lightweight Profile.} Designed for resource-constrained devices such as
IoT sensors and minimal containers. The node is fully stateless: it holds no
persistent local storage. The lifecycle of a lightweight node in each round is:
\emph{register} $\to$ \emph{download global model} $\to$ \emph{train locally}
$\to$ \emph{submit update} $\to$ \emph{disconnect}. On crash, the node simply
re-registers with the coordinator at startup. Recovery is immediate, but the
node cannot resume an interrupted training epoch---any in-progress work is lost.

\textbf{Standard Profile.} Designed for persistent edge appliances and servers
with durable storage. Before transmitting an update, the node writes the
compressed payload, compression algorithm, sample count, and training metrics
to a local SQLite database as a \emph{pending update}. If the node crashes
before the update is acknowledged by the coordinator, the pending record
survives. On restart, the node checks the database for unsent updates and
replays them via the \texttt{SubmitUpdate} RPC, ensuring that training work is
not wasted. Additionally, the standard profile caches the last received global
model checkpoint locally, allowing the node to resume inference immediately
after a restart without waiting for a coordinator broadcast.

The profile is selected via the \texttt{NODE\_PROFILE} environment variable
(defaulting to \texttt{lightweight}), making the choice transparent to the
training and inference logic.

\subsection{Handling Node Dropout and Late Joiners}
\label{subsec:fl_dropout}

Node dropout---caused by network partitions, node crashes, or resource
contention---is a routine occurrence in edge environments. The framework
employs three strategies to maintain round progress.

\textbf{Minimum participation threshold.} The coordinator requires at least
$K_{\min}$ (default~3) client updates before triggering aggregation. If the
threshold is not met within the round timeout (default 120\,s), the round is
aborted and a new round begins. This prevents aggregation from a
non-representative subset of clients.

\textbf{Timeout-bounded waiting.} Rather than blocking indefinitely for all
clients, the coordinator waits up to the round timeout and then aggregates with
whatever updates have been received (provided $K_{\min}$ is met). This
asynchronous-tolerant design means that slow or temporarily disconnected nodes
do not stall the global training loop. Updates that arrive after the round
deadline are cached in Redis (if available) and merged into the next round.

\textbf{Weighted partial aggregation.} Even when only a subset of registered
clients participate, FedAvg aggregation proceeds with sample-count weighting:

\[
    \theta_{t+1} = \sum_{k \in \mathcal{P}_t}
    \frac{n_k}{\sum_{j \in \mathcal{P}_t} n_j} \theta_t^{(k)},
\]

where $\mathcal{P}_t \subseteq \{1, \dots, K\}$ is the set of participating
clients in round~$t$. This ensures that the contribution of each client is
proportional to its dataset size, regardless of how many clients participated.
The scalability benchmarks in Chapter~\ref{ch:evaluation} demonstrate that
aggregation time scales linearly from 5.1\,ms at 5~nodes to 990\,ms at
1\,000~nodes, confirming that partial aggregation remains practical even as the
cluster grows.
